\chapter{Mechanische oscillatoren en het meten van de tijd}\label{ch:g12:oscillators-and-time}

Een staande klok tikt: elke tik is een koperen slinger die de verticaal
kruist, sinds nieuwjaar vijftien miljoen keer herhaald. Het kwartsplaatje in
je horloge buigt \num{32768} keer per seconde; een cesiumatoom zoemt negen
miljard keer sneller. Alles wat slingert kan de tijd bijhouden: dit hoofdstuk
volgt het slingeren, van Huygens tot de \omterm{def:g10:universal-gravitation:satellite}{satellieten}.

\section{Vrije trillingen}

\begin{definition}[Oscillator, periode, frequentie, amplitude]\label{def:g12:oscillators-and-time:oscillator}
Een mechanische \emph{oscillator}\index{oscillator} is een systeem dat, uit
een stabiele \emph{evenwichtsstand}\index{evenwichtsstand} gebracht en
losgelaten, daaromheen heen en weer slingert: dat zijn \emph{vrije
trillingen}\index{vrije trillingen}. De beweging herhaalt zich na de
\emph{periode}\index{periode} $T_0$ (\unit{s}); de
\emph{frequentie}\index{frequentie} is $f_0 = 1/T_0$ (\unit{Hz},
\cref{ch:g12:mechanical-waves}); de \emph{amplitude}\index{amplitude} is de
maximale uitwijking.
\end{definition}

\begin{example}[Een dierentuin van oscillatoren]\label{ex:g12:oscillators-and-time:zoo}
Een schommel op het plein: $T_0 \approx \qty{3}{s}$. De A-snaar van een
gitaar: $f_0 = \qty{110}{Hz}$. Een kwartskristal: \qty{32768}{Hz}. De
cesiumtrilling die de seconde vastlegt: ongeveer \qty{9.2e9}{Hz}. Tien
\omterm{def:g10:orders-of-magnitude:oom}{ordegrootten}, één taak: tellen.
\end{example}

\section{De mathematische slinger}

\begin{definition}[Mathematische slinger]\label{def:g12:oscillators-and-time:pendulum}
Een \emph{mathematische slinger}\index{mathematische slinger}: een klein
lichaam met massa $m$ aan een onrekbare draad met lengte $L$ en
verwaarloosbare massa, slingerend in een verticaal vlak. Van de twee \omterm{def:g10:inertia:force}{krachten}
op het lichaam (\cref{ch:g12:newtons-laws}) staat de \omterm{def:g10:inertia:inventory}{spankracht} loodrecht op
de beweging, terwijl de component van het \omterm{def:g10:universal-gravitation:weight}{gewicht} langs de \omterm{def:g10:relative-motion:trajectory}{baan} terugwijst
naar het laagste punt: het \omterm{def:g10:universal-gravitation:weight}{gewicht} is de \emph{terugdrijvende \omterm{def:g10:inertia:force}{kracht}}%
\index{terugdrijvende kracht}, die de uitwijking tegenwerkt.
\end{definition}

\begin{omfigure}
\begin{tikzpicture}[scale=0.8, font=\small]
  \coordinate (O) at (0,0); \draw (-0.6,0) -- (0.6,0);
  \foreach \x in {-0.5,-0.25,0,0.25,0.5} \draw (\x,0) -- ++(0.14,0.14);
  \fill (O) circle (1.5pt); \draw[dashed] (O) -- (0,-3.4);
  \draw[dashed, omExo] (-120:3.0) arc (-120:-60:3.0);
  \draw (O) -- (-60:3.0) node[midway, above right] {$L$};
  \draw (-90:1.0) arc (-90:-60:1.0); \node at (-75:1.32) {$\theta$};
  \coordinate (B) at (1.5,-2.598); \fill[omDef] (B) circle (2.5pt);
  \draw[-Stealth, omDef, thick] (B) -- ++(-0.75,1.299) node[above left] {$\vect T$};
  \draw[-Stealth, thick] (B) -- ++(0,-1.8) node[below] {$\vect P$};
  \draw[-Stealth, omThm, thick] (B) -- ++(-0.779,-0.45) node[left] {$\vect{P_t}$};
  \draw[dashed, omExo] (0.721,-3.048) -- (1.5,-4.398)
    (B) -- ++(0.779,-1.35) (2.279,-3.948) -- (1.5,-4.398);
\end{tikzpicture}

{\small \omterm{def:g10:inertia:force}{Krachten} op het slingerlichaam: de \omterm{def:g10:inertia:inventory}{spankracht} $\vect T$ staat
loodrecht op de beweging; het tangentiële deel $\vect{P_t}$ van het \omterm{def:g10:universal-gravitation:weight}{gewicht}
wijst altijd terug naar de verticaal --- de \omterm{def:g12:oscillators-and-time:pendulum}{terugdrijvende kracht}.}
\end{omfigure}

\begin{proposition}[Periode van de mathematische slinger]\label{prop:g12:oscillators-and-time:pendulum-period}
\index{mathematische slinger!periode}
Voor kleine \omterm{def:g12:oscillators-and-time:oscillator}{amplitudes} (onder ongeveer $15^\circ$) is de \omterm{def:g12:oscillators-and-time:oscillator}{periode}
\[
T_0 = 2\pi \sqrt{\frac{L}{g}}, \qquad g = \qty{9.81}{m/s^2}.
\]
Ze hangt af van de massa van het slingerlichaam noch --- zolang de hoek klein
blijft --- van de \omterm{def:g12:oscillators-and-time:oscillator}{amplitude}: kleine slingeringen zijn
\emph{isochroon}\index{isochroon}, ze duren allemaal even lang.
\end{proposition}
\admitted

\begin{remark}[Dimensieanalyse, en waar de $2\pi$ woont]\label{rem:g12:oscillators-and-time:dimensional}
De formule laat zich raden. Uit $m$ (\unit{kg}), $L$ (\unit{m}) en $g$
(\unit{m/s^2}) is $\sqrt{L/g}$ de enige combinatie met de \omterm{def:g10:orders-of-magnitude:unit}{eenheid} van een
tijd --- en de massa \emph{kan} er niet in zitten, want geen enkel ander
gegeven draagt een \omterm{def:g10:orders-of-magnitude:unit}{kilogram} om haar weg te delen. Dimensies kunnen het zuivere
getal ervoor niet leveren: de $2\pi$ komt uit het oplossen van de
bewegingsvergelijking, wat in het volume van bachelorjaar 1 gebeurt. En de
voorwaarde van de kleine hoek is niet vrijblijvend: bij $30^\circ$ is de ware
\omterm{def:g12:oscillators-and-time:oscillator}{periode} $2\%$ langer, en is het isochronisme weg.
\end{remark}

\begin{example}[De secondeslinger]\label{ex:g12:oscillators-and-time:seconds}
Welke lengte slaat de seconde, $T_0 = \qty{2.00}{s}$, één tik per kant?
Oplossen geeft $L = g\,T_0^2/4\pi^2 = 9.81 \times 2.00^2/4\pi^2 \approx
\qty{0.994}{m}$. Bijna een \omterm{def:g10:orders-of-magnitude:unit}{meter}: slingerklokken zijn hoog omdat seconden
lang zijn.
\end{example}

\section{De veer met massa}

\begin{definition}[Terugdrijvende kracht van een veer]\label{def:g12:oscillators-and-time:spring}
Een wagentje met massa $m$ glijdt wrijvingsloos over een horizontale rail en
hangt aan een veer met \emph{veerconstante}\index{veerconstante} $k$
(\unit{N/m}). Zij $x$ zijn uitwijking uit het \omterm{def:g10:inertia:compensate}{evenwicht}. Uitgerekt of
ingedrukt trekt of duwt de veer het wagentje terug met de \omterm{def:g10:inertia:force}{kracht} $F = -kx$:
evenredig met de uitwijking en er tegengesteld aan --- opnieuw een
\omterm{def:g12:oscillators-and-time:pendulum}{terugdrijvende kracht}, nu geleverd door de elasticiteit.
\end{definition}

\begin{proposition}[Bewegingsvergelijking en haar oplossing]\label{prop:g12:oscillators-and-time:spring-period}
\index{veeroscillator}
De tweede wet van Newton langs de rail (\cref{ch:g12:newtons-laws}) geeft
$m \frac{d^2x}{dt^2} = -kx$, oftewel $\frac{d^2x}{dt^2} =
-\frac{k}{m}\,x$. Wat de constanten $A$ en $\varphi$ ook zijn,
\[
x(t) = A \cos\!\left(\frac{2\pi t}{T_0} + \varphi\right),
\qquad T_0 = 2\pi\sqrt{\frac{m}{k}},
\]
voldoet aan die vergelijking: een trilling met \omterm{def:g12:oscillators-and-time:oscillator}{periode} $T_0$, onafhankelijk
van de \omterm{def:g12:oscillators-and-time:oscillator}{amplitude}.
\end{proposition}

\begin{proof}
Afleiden geeft $\frac{dx}{dt} = -\frac{2\pi}{T_0}A \sin(\frac{2\pi
t}{T_0} + \varphi)$; nog eens: $\frac{d^2x}{dt^2} = -(\frac{2\pi}{T_0})^2
A\cos(\frac{2\pi t}{T_0} + \varphi) = -(\frac{2\pi}{T_0})^2\,x(t)$, en dat is
$-(k/m)\,x$ precies wanneer $(2\pi/T_0)^2 = k/m$.
\end{proof}

\begin{remark}[Alle bewegingen die er zijn]\label{rem:g12:oscillators-and-time:uniqueness}
Dat \emph{elke} beweging van het wagentje zo'n functie is, wordt hier
toegegeven en in het volume van bachelorjaar 1 bewezen. En geen $g$ in
$2\pi\sqrt{m/k}$: stijver trilt sneller, zwaarder trager, en op de Maan
precies hetzelfde.
\end{remark}

\begin{definition}[Amplitude en fase]\label{def:g12:oscillators-and-time:phase}
In $x(t) = A\cos(2\pi t/T_0 + \varphi)$ is de positieve constante $A$ de
\omterm{def:g12:oscillators-and-time:oscillator}{amplitude} --- de uitersten zijn $x = \pm A$ --- en is $2\pi t/T_0 + \varphi$
de \emph{fase}\index{fase}, die per \omterm{def:g12:oscillators-and-time:oscillator}{periode} met $2\pi$ groeit; de
\emph{beginfase}\index{beginfase} $\varphi$ legt de start vast: uit rust
losgelaten op $x = A$ is $\varphi = 0$. Zo geeft $m = \qty{0.20}{kg}$ aan
$k = \qty{20}{N/m}$ de waarde $T_0 = 2\pi\sqrt{0.20/20} = \qty{0.63}{s}$; tot
\qty{5.0}{cm} uitgetrokken en losgelaten wordt dat $x(t) = 5.0\cos(2\pi
t/0.63)$ in centimeter.
\end{definition}

\section{Energie, demping, resonantie}

\begin{proposition}[Energie van de oscillator]\label{prop:g12:oscillators-and-time:energy}
Een veer die over $x$ is uitgerekt slaat de elastische \omterm{def:g11:mechanical-energy:potential}{potentiële energie}
$E_p = \tfrac12 kx^2$ op (verantwoord in \cref{ch:g12:work-and-energy}).
Tijdens een wrijvingsloze trilling is de \omterm{def:g11:mechanical-energy:em}{mechanische energie}
\[
E_m = E_k + E_p = \tfrac12 m v^2 + \tfrac12 k x^2 = \tfrac12 k A^2
\]
constant: volledig elastisch in de uitersten, volledig kinetisch in het
\omterm{def:g10:inertia:compensate}{evenwicht}. De slinger speelt hetzelfde duet tussen $E_k$ en de
gravitatie-$E_p$ (\cref{ch:g11:mechanical-energy}).
\end{proposition}

\begin{proof}
Met $v = \frac{dx}{dt}$ is de afgeleide van $E_m$ gelijk aan
$mv\frac{dv}{dt} + kx\frac{dx}{dt} = v\,(m\frac{d^2x}{dt^2} + kx) = 0$
wegens de bewegingsvergelijking; de waarde $\tfrac12 kA^2$ lees je in een
uiterste af.
\end{proof}

\begin{omfigure}
\begin{tikzpicture}
\begin{axis}[omaxis, width=8.6cm, height=4.6cm, xlabel={$t/T_0$},
    ylabel={energie $/\,E_m$}, xmin=0, xmax=1, ymin=0, ymax=1.3,
    legend pos=south east]
  \addplot[omDef, thick, domain=0:1, samples=120] {cos(360*x)^2};
  \addplot[omProp, thick, domain=0:1, samples=120] {sin(360*x)^2};
  \addplot[omThm, dashed, thick, domain=0:1, samples=2] {1};
  \legend{$E_p$, $E_k$, $E_m$}
\end{axis}
\end{tikzpicture}

{\small Eén \omterm{def:g12:oscillators-and-time:oscillator}{periode}, losgelaten vanaf $x = A$: de \omterm{def:g10:energy-conservation:energy}{energie} klotst twee keer
van elastisch naar kinetisch en terug, met de som vastgepind op
$\tfrac12 kA^2$.}
\end{omfigure}

\begin{definition}[Gedempte trillingen]\label{def:g12:oscillators-and-time:damping}
\omterm{def:g10:inertia:inventory}{Wrijving} tapt uit elke echte \omterm{def:g12:oscillators-and-time:oscillator}{oscillator} \omterm{def:g11:mechanical-energy:em}{mechanische energie} af: zijn \omterm{def:g12:oscillators-and-time:oscillator}{vrije
trillingen} ondergaan \emph{demping}\index{demping}. Bij zwakke demping trilt
het systeem nog, met een stervende \omterm{def:g12:oscillators-and-time:oscillator}{amplitude} maar een nauwelijks veranderde
herhaaltijd, de \emph{pseudoperiode}\index{pseudoperiode}, dicht bij $T_0$;
bij sterke demping (honing) kruipt het naar het \omterm{def:g10:inertia:compensate}{evenwicht} terug zonder er
ooit voorbij te schieten; de \omterm{def:g10:energy-conservation:energy}{energie}, evenredig met het kwadraat van de
\omterm{def:g12:oscillators-and-time:oscillator}{amplitude}, vervalt tot warmte.
\end{definition}

\begin{omfigure}
\begin{tikzpicture}
\begin{axis}[omaxis, width=8.6cm, height=4.6cm, xlabel={$t$ (\unit{s})},
    ylabel={$x$ (\unit{cm})}, xmin=0, xmax=10, ymin=-2.6, ymax=2.6,
    legend pos=south east]
  \addplot[omExo, thin, domain=0:10, samples=240] {2*cos(deg(pi*x))};
  \addplot[omDef, thick, domain=0:10, samples=240] {2*exp(-x/4)*cos(deg(pi*x))};
  \legend{ongedempt, gedempt}
  \addplot[omThm, dashed, domain=0:10, samples=60, forget plot] {2*exp(-x/4)};
  \addplot[omThm, dashed, domain=0:10, samples=60, forget plot] {-2*exp(-x/4)};
\end{axis}
\end{tikzpicture}

{\small Trillingen zonder en met zwakke \omterm{def:g12:oscillators-and-time:damping}{demping}: vrijwel dezelfde \omterm{def:g12:oscillators-and-time:oscillator}{periode}
(hier \qty{2.0}{s}), met de gedempte \omterm{def:g12:oscillators-and-time:oscillator}{amplitude} die binnen een krimpende
exponentiële omhullende (gestippeld) sterft.}
\end{omfigure}

\begin{definition}[Gedwongen trillingen en resonantie]\label{def:g12:oscillators-and-time:resonance}
Duw een \omterm{def:g12:oscillators-and-time:oscillator}{oscillator} \omterm{def:g10:signals-and-waves:period}{periodiek} aan --- \emph{gedwongen
trillingen}\index{gedwongen trillingen} --- en na een korte inloop trilt hij
met de \emph{aandrijvende} \omterm{def:g12:oscillators-and-time:oscillator}{frequentie} $f$, niet met de zijne. Zijn \omterm{def:g12:oscillators-and-time:oscillator}{amplitude}
piekt wanneer $f$ de \emph{eigenfrequentie}\index{eigenfrequentie} $f_0$ van
de vrije \omterm{def:g12:oscillators-and-time:oscillator}{oscillator} nadert: dat is \emph{resonantie}\index{resonantie},
waarbij elke duw in de pas met de beweging aankomt en elke cyclus \omterm{def:g10:energy-conservation:energy}{energie}
toevoert; hoe zwakker de \omterm{def:g12:oscillators-and-time:damping}{demping}, hoe hoger en scherper de piek.
\end{definition}

\begin{omfigure}
\begin{tikzpicture}
\begin{axis}[omaxis, width=8.6cm, height=4.6cm, xlabel={$f/f_0$},
    ylabel={amplitude (willek.)}, xmin=0, xmax=2.2, ymin=0, ymax=5.6,
    legend pos=north east]
  \addplot[omThm, thick, domain=0:2.2, samples=200] {1/sqrt((1-x^2)^2 + (x/5)^2)};
  \addplot[omDef, thick, domain=0:2.2, samples=120] {1/sqrt((1-x^2)^2 + (x/1.2)^2)};
  \legend{zwakke demping, sterke demping}
  \addplot[omExo, dashed, forget plot] coordinates {(1,0) (1,5.4)};
\end{axis}
\end{tikzpicture}

{\small Resonantiekromme: de gedwongen \omterm{def:g12:oscillators-and-time:oscillator}{amplitude} tegen de aandrijvende
\omterm{def:g12:oscillators-and-time:oscillator}{frequentie}. Bij $f = f_0$ torent het antwoord omhoog, hoger naarmate de
\omterm{def:g12:oscillators-and-time:damping}{demping} zwakker is.}
\end{omfigure}

\begin{example}[Resonantie ten goede en ten kwade]\label{ex:g12:oscillators-and-time:resonance-zoo}
Een ouder die een schommel aanduwt, duwt op diens eigen \omterm{def:g12:oscillators-and-time:oscillator}{frequentie} ---
\omterm{def:g12:oscillators-and-time:resonance}{resonantie} aan het werk. In \num{2000} zwaaide de nieuwe Millennium Bridge in
Londen verontrustend toen de passen van de menigte op een \omterm{def:g12:oscillators-and-time:resonance}{eigenfrequentie}
rond \qty{1}{Hz} inhaakten; hij sloot twee jaar voor dempers. Elk
kwartshorloge houdt zijn kristal aan het zingen door het op \omterm{def:g12:oscillators-and-time:resonance}{resonantie} aan te
drijven, \qty{32768}{Hz}.
\end{example}

\section{De tijd meten}

\begin{method}[Een periode opmeten]\label{met:g12:oscillators-and-time:timing}
Meet nooit één enkele trilling: start de chronometer wanneer het
slingerlichaam de verticaal kruist (daar is het het snelst en het best te
beoordelen), tel $N = 50$ trillingen en deel door $N$ --- een reactiefout van
\qty{0.2}{s} wordt zo \qty{4}{ms} op $T_0$. Herhaal, en vergelijk de
metingen.
\end{method}

\begin{example}[Drie eeuwen tik]\label{ex:g12:oscillators-and-time:clocks}
Huygens bouwde in \num{1657} de eerste slingerklok: het isochronisme bracht
klokken van een kwartier afwijking per dag naar seconden, en gecompenseerde
slingers naar fracties van een seconde per dag. Het kwartspolshorloge
(\num{1969}) verving de staaf door een kristal dat met \qty{32768}{Hz} buigt:
tienden van een seconde per maand. Sinds \num{1967} wordt de seconde zelf
\emph{gedefinieerd} door een atomaire trilling: de duur van \num{9192631770}
\omterm{def:g12:oscillators-and-time:oscillator}{perioden} van de microgolfstraling van een inwendige overgang van cesium-133
(\cref{ch:g12:quantum-world}). De beste cesiumklokken lopen één seconde per
honderd miljoen jaar uit de pas.
\end{example}

\begin{remark}[Waarom de nanoseconde najagen]\label{rem:g12:oscillators-and-time:gps}
Een gps-ontvanger bepaalt zijn plaats door radiosignalen van \omterm{def:g10:universal-gravitation:satellite}{satellieten} te
klokken; licht legt \qty{30}{cm} per nanoseconde af, dus zet een klokfout van
één microseconde je \qty{300}{m} verkeerd. \omterm{def:g10:universal-gravitation:satellite}{Satellieten} dragen daarom
atoomklokken --- en bij die nauwkeurigheid hangt het tempo van een klok af van
hoe snel ze beweegt en hoe hoog ze zit, een effect dat we in
\cref{ch:g12:special-relativity} tegenkomen.
\end{remark}

\section{Oefeningen}

\begin{exercise}[$\star$]\label{exo:g12:oscillators-and-time:1}
(a) Bereken de \omterm{def:g12:oscillators-and-time:oscillator}{periode} en de \omterm{def:g12:oscillators-and-time:oscillator}{frequentie} van een \omterm{def:g12:oscillators-and-time:pendulum}{mathematische slinger} van
\qty{1.00}{m}. (b) De lengte wordt verviervoudigd: wat gebeurt er met de
\omterm{def:g12:oscillators-and-time:oscillator}{periode}?
\end{exercise}

\begin{exercise}[$\star$]\label{exo:g12:oscillators-and-time:2}
Een hart in rust slaat $72$ keer per minuut; de vleugel van een libel met
\qty{30}{Hz}. Geef de \omterm{def:g12:oscillators-and-time:oscillator}{frequentie} en de \omterm{def:g12:oscillators-and-time:oscillator}{periode} van het hart en de \omterm{def:g12:oscillators-and-time:oscillator}{periode} van
de vleugel; welk woord uit dit hoofdstuk benoemt de maximale rijzing van de
borstkas per slag?
\end{exercise}

\begin{exercise}[$\star$]\label{exo:g12:oscillators-and-time:3}
Een wagentje van \qty{0.10}{kg} hangt aan een veer van \qty{25}{N/m}: bereken
$T_0$ en $f_0$. Wat doet het verdubbelen van de massa met de \omterm{def:g12:oscillators-and-time:oscillator}{periode}?
\end{exercise}

\begin{exercise}[$\star$]\label{exo:g12:oscillators-and-time:4}
Opeenvolgende maxima van een gedempte \omterm{def:g12:oscillators-and-time:oscillator}{oscillator}: $2.0$, $1.5$, $1.13$,
\qty{0.85}{cm} op $t = 0$, $2.0$, $4.0$, \qty{6.0}{s}. Lees de \omterm{def:g12:oscillators-and-time:damping}{pseudoperiode}
af; toon aan dat de \omterm{def:g12:oscillators-and-time:oscillator}{amplitude} per \omterm{def:g12:oscillators-and-time:oscillator}{periode} met een vaste factor daalt;
voorspel het volgende maximum.
\end{exercise}

\begin{exercise}[$\star$]\label{exo:g12:oscillators-and-time:5}
Een astronaut neemt een slinger en een \omterm{prop:g12:oscillators-and-time:spring-period}{veeroscillator} mee naar de Maan
($g = \qty{1.62}{m/s^2}$). Met welke factor verandert elke \omterm{def:g12:oscillators-and-time:oscillator}{periode}? Verklaar.
\end{exercise}

\begin{exercise}[$\star\star$]\label{exo:g12:oscillators-and-time:6}
(a) Ga na dat $\sqrt{L/g}$ de dimensie van een tijd heeft. (b) Waarom kan
geen enkele formule voor $T_0$ die uit $m$, $L$ en $g$ is opgebouwd de massa
bevatten? (c) Zou dit argument ooit de $2\pi$ kunnen opleveren?
\end{exercise}

\begin{exercise}[$\star\star$]\label{exo:g12:oscillators-and-time:7}
Ga door twee keer af te leiden na dat $x(t) = A\cos(2\pi t/T_0 + \varphi)$
aan $\frac{d^2x}{dt^2} = -(k/m)\,x$ voldoet dan en slechts dan als
$T_0 = 2\pi\sqrt{m/k}$.
\end{exercise}

\begin{exercise}[$\star\star$]\label{exo:g12:oscillators-and-time:8}
Een veer met \omterm{def:g12:oscillators-and-time:spring}{veerconstante} \qty{40}{N/m} draagt een wagentje van
\qty{0.20}{kg} met \omterm{def:g12:oscillators-and-time:oscillator}{amplitude} \qty{3.0}{cm}. Bereken de \omterm{def:g11:mechanical-energy:em}{mechanische energie} en
de maximale snelheid. Waar is de snelheid maximaal? En nul?
\end{exercise}

\begin{exercise}[$\star\star$]\label{exo:g12:oscillators-and-time:9}
(a) Bereken de lengte van de secondeslinger ($T_0 = \qty{2.00}{s}$). (b) Haar
lengte groeit met $1.0\%$: bepaal met $\sqrt{1+x} \approx 1 + x/2$ de
relatieve verandering van $T_0$ en de fout per dag.
\end{exercise}

\begin{exercise}[$\star\star$]\label{exo:g12:oscillators-and-time:10}
Een zwak gedempte \omterm{def:g12:oscillators-and-time:oscillator}{oscillator} verliest per \omterm{def:g12:oscillators-and-time:oscillator}{periode} $5.0\%$ van zijn \omterm{def:g10:energy-conservation:energy}{energie}.
(a) Hoeveel daalt de \omterm{def:g12:oscillators-and-time:oscillator}{amplitude} per \omterm{def:g12:oscillators-and-time:oscillator}{periode}? (b) Na hoeveel \omterm{def:g12:oscillators-and-time:oscillator}{perioden} is de
helft van de \omterm{def:g10:energy-conservation:energy}{energie} weg? (c) Ligt de \omterm{def:g12:oscillators-and-time:damping}{pseudoperiode} ver van $T_0$?
\end{exercise}

\begin{exercise}[$\star\star$]\label{exo:g12:oscillators-and-time:11}
Verklaar met de resonantiekromme: (a) waarom een schommel aanduwen alleen op
haar eigen ritme werkt; (b) waarom soldaten op een loopbrug uit de pas gaan
lopen; (c) waarom een wasmachine bij één toerental staat te schudden terwijl
ze op gang komt.
\end{exercise}

\begin{exercise}[$\star\star\star$]\label{exo:g12:oscillators-and-time:12}
Een slingerklok die op zeeniveau juist loopt, verhuist naar een sterrenwacht
waar $g = \qty{9.78}{m/s^2}$. (a) Loopt ze voor of achter? Hoeveel seconden
per dag? (b) Hoeveel korter moet haar slinger van \qty{0.994}{m} worden?
\end{exercise}

\begin{exercise}[$\star\star\star$]\label{exo:g12:oscillators-and-time:13}
Een wagentje van \qty{0.20}{kg} voldoet aan $x(t) = A\cos(2\pi t/T_0)$, met
$A = \qty{4.0}{cm}$ en $T_0 = \qty{0.63}{s}$. Bereken de maximale snelheid,
de maximale versnelling, de \omterm{def:g12:oscillators-and-time:spring}{veerconstante} $k$ en de \omterm{def:g11:mechanical-energy:em}{mechanische energie}.
\end{exercise}

\begin{exercise}[$\star\star\star$]\label{exo:g12:oscillators-and-time:14}
Om $g$ te meten: een slinger met $L = 99.5 \pm \qty{0.2}{cm}$ maakt $20$
trillingen in $40.1 \pm \qty{0.2}{s}$. Bereken $g$ en de \omterm{def:g10:orders-of-magnitude:uncertainty}{relatieve
onzekerheid} ervan (tel die van $L$ en die van $T_0^2$ op); is het resultaat
in overeenstemming met \qty{9.81}{m/s^2}?
\end{exercise}

\begin{exercise}[$\star\star\star$]\label{exo:g12:oscillators-and-time:15}
Vergelijk tijdmeters via hun relatieve afwijking (de fout gedeeld door de
verstreken tijd): een slingerklok die \qty{10}{s} per dag verliest, een
kwartshorloge \qty{0.5}{s} per maand, een cesiumklok \qty{1}{s} per honderd
miljoen jaar: bereken de drie verhoudingen. Een gps-klok die er
\qty{1}{\micro s} naast zit: hoe groot is de plaatsfout? Welke familie heeft
gps nodig?
\end{exercise}

\section{Probleem: De werkbank van de klokkenmaker}

\begin{problem}[{Weekendopgave --- de werkbank van de klokkenmaker: een
wandklok weer op tijd gebracht --- vijftig slingeringen geklokt, een koperen
staaf die in de zomer uitzet, het gefluisterde duwtje van het echappement,
een uitdager van kwarts --- en het oordeel, uitgesproken in seconden per
maand}]
\label{pb:g12:oscillators-and-time:1}
Een slingerwandklok komt ``verkeerd lopend'' op de werkbank. Behandel haar
koperen staaf en zware slingerlichaam als een \omterm{def:g12:oscillators-and-time:pendulum}{mathematische slinger} met
effectieve lengte $L$ en massa $m = \qty{1.2}{kg}$; het raderwerk telt elke
volledige trilling als precies \qty{2.000}{s}.

\medskip
\noindent\textbf{Deel I --- De klok de pols voelen.}
\begin{enumerate}
  \item Waarom $50$ trillingen klokken en niet één? Schat de fout die op $T_0$
        overblijft als je reactietijd \qty{0.2}{s} is.
  \item De chronometer wijst \qty{99.4}{s} voor $50$ trillingen: hoe groot is
        de \omterm{def:g12:oscillators-and-time:oscillator}{periode}?
  \item Leid daaruit de huidige effectieve lengte van de slinger af.
  \item Welke lengte zou precies \qty{2.000}{s} geven? Moet de stelmoer het
        slingerlichaam omhoog of omlaag brengen, en met hoeveel?
  \item Loopt de klok, onbijgesteld, voor of achter, en met hoeveel seconden
        per dag?
  \item \omterm{def:g10:signals-and-waves:oscilloscope}{Gevoeligheid} voor $g$: wat is de \omterm{def:g12:oscillators-and-time:oscillator}{periode} op de Maan
        ($g = \qty{1.62}{m/s^2}$)? En de fout per dag op een sterrenwacht waar
        $g = \qty{9.80}{m/s^2}$?
\end{enumerate}

\medskip
\noindent\textbf{Deel II --- De zomerse inzinking.} Koper zet uit: een
opwarming $\Delta\theta$ rekt de staaf tot $L' = L\,(1 + \lambda\,
\Delta\theta)$, met $\lambda = \num{1.9e-5}$ per \omterm{def:g10:light-spectra:kelvin}{kelvin}. De klok is bij
\qty{20}{\celsius} bijgesteld.
\begin{enumerate}[resume]
  \item Bereken $\Delta L$ voor een zomerkamer op \qty{25}{\celsius}.
  \item Toon met $\sqrt{1+x} \approx 1 + x/2$ aan dat
        $\Delta T_0 / T_0 = \tfrac12 \lambda\, \Delta\theta$.
  \item Bereken $\Delta T_0/T_0$ en de zomerse fout per dag: voor of achter?
  \item Dezelfde vragen voor een wintergang op \qty{10}{\celsius}.
  \item Oude precisieklokken gebruikten ``roosterslingers'' die koperen staven
        met stalen (die ongeveer half zoveel uitzetten) in tegengestelde
        richtingen combineerden. Leg uit hoe dat het tempo kan vasthouden.
\end{enumerate}

\medskip
\noindent\textbf{Deel III --- Het gefluisterde duwtje.} Met het echappement
ontkoppeld halveert de slingering, gestart met \omterm{def:g12:oscillators-and-time:oscillator}{amplitude}
$\theta_m = 2.00^\circ$, haar \omterm{def:g12:oscillators-and-time:oscillator}{amplitude} in \qty{280}{s}; in bedrijf houdt een
klein duwtje van het echappement per \omterm{def:g12:oscillators-and-time:oscillator}{periode} $\theta_m$ constant.
\begin{enumerate}[resume]
  \item Hoeveel \omterm{def:g12:oscillators-and-time:oscillator}{perioden} duurt die halvering?
  \item In het uiterste is het slingerlichaam $h = L(1 - \cos\theta_m)
        \approx L\,\theta_m^2/2$ geklommen (in radiaal): leid
        $E_m = \tfrac12 m g L\,\theta_m^2$ af en bereken die
        (\cref{ch:g11:mechanical-energy}).
  \item De \omterm{def:g10:energy-conservation:energy}{energie} is evenredig met het kwadraat van de \omterm{def:g12:oscillators-and-time:oscillator}{amplitude}: welke
        fractie van $E_m$ gaat er per \omterm{def:g12:oscillators-and-time:oscillator}{periode} verloren?
  \item Leid de \omterm{def:g10:energy-conservation:energy}{energie} van elk duwtje af, en het gemiddelde \omterm{def:g11:work-of-force:power}{vermogen} dat
        nodig is.
  \item Het aandrijfgewicht $M = \qty{2.5}{kg}$ zakt $h = \qty{1.0}{m}$ per
        week: is dat genoeg \omterm{def:g10:energy-conservation:energy}{energie}, en bij welk \omterm{def:g10:energy-conservation:efficiency}{rendement}?
\end{enumerate}

\medskip
\noindent\textbf{Deel IV --- De uitdager van kwarts.} Een kwartsuurwerk trilt
met $f = \qty{32768}{Hz}$ en wijkt ongeveer \qty{0.3}{s} per maand af; de
herstelde slingerklok wijkt nog altijd \qty{0.5}{s} per dag af.
\begin{enumerate}[resume]
  \item Bereken de \omterm{def:g12:oscillators-and-time:oscillator}{periode} van het kwarts. Toon aan dat $\num{32768} = 2^{15}$:
        waarom is een macht van twee precies wat een horlogeschakeling wil?
  \item Bereken de relatieve afwijkingen van het kwarts en van de herstelde
        slingerklok, en beide in seconden per maand. Wie wint, en met welke
        factor?
  \item De cesiumstandaard --- \num{9192631770} \omterm{def:g12:oscillators-and-time:oscillator}{perioden} maken één seconde ---
        wijkt \qty{1}{s} per honderd miljoen jaar af: hoe groot is haar
        relatieve afwijking? Hoeveel \omterm{def:g10:orders-of-magnitude:oom}{ordegrootten} onder het kwarts, en waarom
        eist gps haar?
  \item Het oordeel, in één zin met getallen: wat verdient de herstelde klok
        aan de muur, en van wie is de seconde nu?
\end{enumerate}
\end{problem}
